% GENERATED FILE. Edit canonical lessons, then run npm run generate:books.
% canonical-source-hash: fnv1a64:24d29eed35707ddc
% canonical-lessons: UR-C454-fitrat, UR-C454-tabiat, UR-C454-adat, UR-C454-khubi, UR-C454-shakhsiyat

\chapter{Character}
\label{ch:ur-character}
\hlchaptermodality{454}

\begin{chapteropening}
\textbf{By the end of this chapter:} \emph{I can describe someone's character.}

Complete the last lesson of chapter 454: I can describe someone's character.
\end{chapteropening}

\section[\texorpdfstring{\ur{فطرت} (fitrat)}{فطرت (fitrat)}]{\ur{فطرت} (fitrat) — a person's nature; disposition}

\label{lesson:UR-C454-fitrat}



\begin{quote}
Before the new one: say the Urdu for to sob, then the Urdu for teasing.
\end{quote}

\subsection*{You'll want to know: \ur{فطرت}}
\textbf{\ur{فطرت}} — \emph{fitrat} — \textquotedblleft{}a person's nature; disposition\textquotedblright{}.

\begin{quote}
\ur{اس} \ur{کی} \ur{فطرت} \ur{بڑی} \ur{نرم} \ur{ہے۔}
\end{quote}

\begin{quote}
\emph{us kī fitrat baṛī narm hai.}
\end{quote}

\begin{quote}
She has a very gentle nature.
\end{quote}

Feminine.

\begin{cousinweb}
\textbf{\ur{فطرت}} (\emph{fitrat}) is usually traced, through Persian, to Arabic \emph{fiṭra}, \textquotedblleft{}the nature a person is born with\textquotedblright{}.
\end{cousinweb}

\begin{grammarlens}[title={What you've built}]
One more word for describing what you lived and felt.
\end{grammarlens}

\subsection*{Your turn}
\begin{itemize}
\raggedright
  \item \emph{Say it:} \emph{fitrat}
  \item \emph{Say it:} \emph{fitrat}, once more
  \item \emph{Say it:} \emph{us kī fitrat baṛī narm hai.} — She has a very gentle nature.
  \item \emph{Recall:} say the Urdu for to sob, then the Urdu for teasing, then say \emph{fitrat} again
\end{itemize}

\begin{tcolorbox}[breakable,colback=teal!4,colframe=teal!35!black,title={Before you move on}]
What does \ur{فطرت} mean? (\textquotedblleft{}A person's nature; disposition\textquotedblright{}.) Say it once more.
\end{tcolorbox}
\section[\texorpdfstring{\ur{طبیعت} (tabīat)}{طبیعت (tabīat)}]{\ur{طبیعت} (tabīat) — disposition; how one feels (in health)}

\label{lesson:UR-C454-tabiat}



\begin{quote}
Before the new one: say the Urdu for teasing, then the Urdu for a person's nature.
\end{quote}

\subsection*{You'll want to know: \ur{طبیعت}}
\textbf{\ur{طبیعت}} — \emph{tabīat} — \textquotedblleft{}disposition; how one feels (in health)\textquotedblright{}.

\begin{quote}
\ur{آج} \ur{میری} \ur{طبیعت} \ur{ٹھیک} \ur{نہیں} \ur{ہے۔}
\end{quote}

\begin{quote}
\emph{āj merī tabīat ṭhīk nahīṅ hai.}
\end{quote}

\begin{quote}
I'm not feeling well today.
\end{quote}

Feminine.

\begin{cousinweb}
\textbf{\ur{طبیعت}} (\emph{tabīat}) is usually traced, through Persian, to Arabic \emph{ṭabīʿa}, \textquotedblleft{}nature, temperament\textquotedblright{}; Urdu also uses it to ask how someone feels.
\end{cousinweb}

\begin{grammarlens}[title={What you've built}]
One more word for describing what you lived and felt.
\end{grammarlens}

\subsection*{Your turn}
\begin{itemize}
\raggedright
  \item \emph{Say it:} \emph{tabīat}
  \item \emph{Say it:} \emph{tabīat}, once more
  \item \emph{Say it:} \emph{āj merī tabīat ṭhīk nahīṅ hai.} — I'm not feeling well today.
  \item \emph{Recall:} say the Urdu for teasing, then the Urdu for a person's nature, then say \emph{tabīat} again
\end{itemize}

\begin{tcolorbox}[breakable,colback=teal!4,colframe=teal!35!black,title={Before you move on}]
What does \ur{طبیعت} mean? (\textquotedblleft{}Disposition; how one feels (in health)\textquotedblright{}.) Say it once more.
\end{tcolorbox}
\section[\texorpdfstring{\ur{عادت} (ādat)}{عادت (ādat)}]{\ur{عادت} (ādat) — a habit}

\label{lesson:UR-C454-adat}



\begin{quote}
Before the new one: say the Urdu for a person's nature, then the Urdu for disposition.
\end{quote}

\subsection*{You'll want to know: \ur{عادت}}
\textbf{\ur{عادت}} — \emph{ādat} — \textquotedblleft{}a habit\textquotedblright{}.

\begin{quote}
\ur{اسے} \ur{صبح} \ur{سویرے} \ur{اٹھنے} \ur{کی} \ur{عادت} \ur{ہے۔}
\end{quote}

\begin{quote}
\emph{use subah savere uṭhne kī ādat hai.}
\end{quote}

\begin{quote}
She is in the habit of getting up early.
\end{quote}

Feminine.

\begin{cousinweb}
\textbf{\ur{عادت}} (\emph{ādat}) is usually traced, through Persian, to Arabic \emph{ʿāda}, \textquotedblleft{}custom, habit\textquotedblright{}.
\end{cousinweb}

\begin{grammarlens}[title={What you've built}]
One more word for describing what you lived and felt.
\end{grammarlens}

\subsection*{Your turn}
\begin{itemize}
\raggedright
  \item \emph{Say it:} \emph{ādat}
  \item \emph{Say it:} \emph{ādat}, once more
  \item \emph{Say it:} \emph{use subah savere uṭhne kī ādat hai.} — She is in the habit of getting up early.
  \item \emph{Recall:} say the Urdu for a person's nature, then the Urdu for disposition, then say \emph{ādat} again
\end{itemize}

\begin{tcolorbox}[breakable,colback=teal!4,colframe=teal!35!black,title={Before you move on}]
What does \ur{عادت} mean? (\textquotedblleft{}A habit\textquotedblright{}.) Say it once more.
\end{tcolorbox}
\section[\texorpdfstring{\ur{خوبی} (khūbī)}{خوبی (khūbī)}]{\ur{خوبی} (khūbī) — a good quality; a strength}

\label{lesson:UR-C454-khubi}



\begin{quote}
Before the new one: say the Urdu for disposition, then the Urdu for a habit.
\end{quote}

\subsection*{You'll want to know: \ur{خوبی}}
\textbf{\ur{خوبی}} — \emph{khūbī} — \textquotedblleft{}a good quality; a strength\textquotedblright{}.

\begin{quote}
\ur{سچ} \ur{بولنا} \ur{اس} \ur{کی} \ur{سب} \ur{سے} \ur{بڑی} \ur{خوبی} \ur{ہے۔}
\end{quote}

\begin{quote}
\emph{sach bolnā us kī sab se baṛī khūbī hai.}
\end{quote}

\begin{quote}
Telling the truth is his greatest strength.
\end{quote}

Feminine.

\begin{cousinweb}
\textbf{\ur{خوبی}} (\emph{khūbī}) is usually traced to Persian \emph{khūbī}, \textquotedblleft{}goodness\textquotedblright{}, built on \emph{khūb}, \textquotedblleft{}good, fine\textquotedblright{}.
\end{cousinweb}

\begin{grammarlens}[title={What you've built}]
One more word for describing what you lived and felt.
\end{grammarlens}

\subsection*{Your turn}
\begin{itemize}
\raggedright
  \item \emph{Say it:} \emph{khūbī}
  \item \emph{Say it:} \emph{khūbī}, once more
  \item \emph{Say it:} \emph{sach bolnā us kī sab se baṛī khūbī hai.} — Telling the truth is his greatest strength.
  \item \emph{Recall:} say the Urdu for disposition, then the Urdu for a habit, then say \emph{khūbī} again
\end{itemize}

\begin{tcolorbox}[breakable,colback=teal!4,colframe=teal!35!black,title={Before you move on}]
What does \ur{خوبی} mean? (\textquotedblleft{}A good quality; a strength\textquotedblright{}.) Say it once more.
\end{tcolorbox}
\section[\texorpdfstring{\ur{شخصیت} (shakhsiyat)}{شخصیت (shakhsiyat)}]{\ur{شخصیت} (shakhsiyat) — a personality}

\label{lesson:UR-C454-shakhsiyat}



\begin{quote}
Before the new one: say the Urdu for a habit, then the Urdu for a good quality.
\end{quote}

\subsection*{You'll want to know: \ur{شخصیت}}
\textbf{\ur{شخصیت}} — \emph{shakhsiyat} — \textquotedblleft{}a personality\textquotedblright{}.

\begin{quote}
\ur{ہمارے} \ur{استاد} \ur{کی} \ur{شخصیت} \ur{بڑی} \ur{دلچسپ} \ur{ہے۔}
\end{quote}

\begin{quote}
\emph{hamāre ustād kī shakhsiyat baṛī dilchasp hai.}
\end{quote}

\begin{quote}
Our teacher has a very interesting personality.
\end{quote}

Feminine.

\begin{cousinweb}
\textbf{\ur{شخصیت}} (\emph{shakhsiyat}) is usually traced, through Persian, to Arabic \emph{shakhṣiyya}, \textquotedblleft{}personality\textquotedblright{}, built on \emph{shakhṣ}, \textquotedblleft{}a person\textquotedblright{}.
\end{cousinweb}

\begin{grammarlens}[title={What you've built}]
One more word for describing what you lived and felt.
\end{grammarlens}

\subsection*{Your turn}
\begin{itemize}
\raggedright
  \item \emph{Say it:} \emph{shakhsiyat}
  \item \emph{Say it:} \emph{shakhsiyat}, once more
  \item \emph{Say it:} \emph{hamāre ustād kī shakhsiyat baṛī dilchasp hai.} — Our teacher has a very interesting personality.
  \item \emph{Recall:} say the Urdu for a habit, then the Urdu for a good quality, then say \emph{shakhsiyat} again
\end{itemize}

\begin{tcolorbox}[breakable,colback=teal!4,colframe=teal!35!black,title={Before you move on}]
What does \ur{شخصیت} mean? (\textquotedblleft{}A personality\textquotedblright{}.) Say it once more.
\end{tcolorbox}
